\section{An elementary projective degree bound for alternating quotients}
\label{sec:projective-abelian-rank}

\paragraph{Author freeze and origin.}
D-TD32, reserved question Q2ep, is a derivation by \texttt{/root}.
It supplies the projective degree input needed for the proposed
homogeneous multiplicity branch. The complete author proof reading
and the independent whole review pass. No finite mathematical
computation is needed or claimed.
The proof is elementary and uses no classification of modular
irreducible representations or Schur multipliers.

\begin{lemma}[D-TD32]
Let $F$ be an algebraically closed field, let $\ell$ be an odd
prime different from $\operatorname{char}F$, and let
\[
                  E\simeq(C_\ell)^k\ \hookrightarrow
                         \operatorname{PGL}_r(F).
 \tag{TDA.1}
\]
Then
\[
                                  k\le r-1.
 \tag{TDA.2}
\]
Consequently, for every $n\ge5$, every faithful projective
representation of $A_n$ over any field of characteristic $p\ge0$
has degree
\[
 r\ge
 \begin{cases}
 \lfloor n/3\rfloor+1,&p\ne3,\\
 \lfloor n/5\rfloor+1,&p=3,
 \end{cases}
 \qquad\text{and in particular}\qquad r>n/5.
 \tag{TDA.3}
\]
The representation need not be irreducible.
\end{lemma}

\begin{proof}
Choose a basis of $E$ as an $\mathbf F_\ell$-vector space and
matrix lifts of its basis vectors. Since each basis vector has
projective order $\ell$, scale its lift to have $\ell$th power
one. This is possible over $F$. Any two lifts commute up to a
scalar, and their commutator scalar is an $\ell$th root of unity:
if their lifts are $X,Y$, then
$[X^\ell,Y]=[X,Y]^\ell=1$. The scalar commutators therefore define
an alternating bilinear pairing
\[
 \beta:E\times E\longrightarrow\mu_\ell(F),\qquad
 \dim_{\mathbf F_\ell}\operatorname{rad}\beta=k-2s.
 \tag{TDA.4}
\]
Identify $\mu_\ell(F)$ with the additive group of $\mathbf F_\ell$
by choosing a primitive $\ell$th root. The even rank $2s$ and a
symplectic basis of a complement to the radical follow by the
elementary alternating-form reduction.

For clarity, the normalized lifts generate a finite group:
every product can be reordered into one of at most
$\ell^{k+1}$ words consisting of a scalar $\ell$th root followed
by the basis lifts with exponents between zero and $\ell-1$.
All its elements have exponent dividing $\ell$. Indeed in this
group of nilpotence class at most two, the scalar correction in
the $\ell$th power of a product has exponent divisible by
$\binom{\ell}{2}$, which is divisible by $\ell$ because $\ell$
is odd. This also permits order-$\ell$ lifts of any new basis.

Put $R_0=\operatorname{rad}\beta$ and choose order-$\ell$ lifts
of its basis. They commute with one another and with every
basis lift of $E$. They therefore give an ordinary linear
representation of $R_0$, central in the lifted group. Since
$\operatorname{char}F\ne\ell$, these matrices are simultaneously
diagonalizable. Decompose the whole vector space according to
its distinct occurring $R_0$-characters:
\[
                 F^r=\bigoplus_{i=1}^t V_{\chi_i},
                       \qquad t\ge1.
 \tag{TDA.5}
\]
Every $V_{\chi_i}$ is stable under all the lifts.

Take a symplectic basis of the $2s$-dimensional complement
to $R_0$, with normalized lifts
$X_1,\ldots,X_s,Z_1,\ldots,Z_s$ such that the $X_j$ commute
with one another, the $Z_j$ commute with one another, and
$[X_i,Z_j]=\zeta^{\delta_{ij}}I$.
On any nonzero $V_{\chi_i}$, simultaneous diagonalization of
the $Z_j$ gives a nonzero common weight space. The $X_j$ translate
its weight independently in all $s$ coordinates. Thus all
$\ell^s$ weights occur, with equal positive multiplicity.
It follows directly that
\[
                \dim V_{\chi_i}=a m_i,\qquad
                a=\ell^s,\qquad m_i\ge1,\qquad r\ge at.
 \tag{TDA.6}
\]
This argument does not assume that the full $E$-representation
is irreducible and does not use a modular classification.

Projective faithfulness is essential for the next count. If a
radical element has the same value under all the $\chi_i$,
its lift is a scalar on all of $F^r$, so its image in
$\operatorname{PGL}_r(F)$ is the identity. By (TDA.1) that
element is zero. The character-ratio homomorphism is therefore
injective:
\[
 R_0\longrightarrow\mu_\ell(F)^{\,t-1},\qquad
 v\longmapsto
   \bigl(\chi_2(v)/\chi_1(v),\ldots,\chi_t(v)/\chi_1(v)\bigr).
 \tag{TDA.7}
\]
Hence $\dim R_0\le t-1$.
The elementary inequality $\ell^s\ge2s+1$ for odd $\ell$
follows by induction from $s=0$: multiplication by at least
three carries $2s+1$ to at least $2(s+1)+1$. Consequently
\[
 k=2s+\dim R_0
       \le2s+t-1
       \le(2s+1)t-1
       \le at-1
       \le r-1,
 \tag{TDA.8}
\]
which proves (TDA.2), including the radical-free and totally
commuting cases.

For the alternating conclusion, extend the coefficient field to
an algebraic closure. This does not change projective faithfulness:
an original matrix scalar after extension was already scalar
over the original field, as any diagonal entry shows.
If $p\ne3$, the subgroup generated by $\lfloor n/3\rfloor$
disjoint three-cycles is elementary abelian of that rank and
lies in $A_n$. Apply (TDA.2) with $\ell=3$.
If $p=3$, use instead $\lfloor n/5\rfloor$ disjoint five-cycles
and $\ell=5$. Odd cycles are even permutations in both cases.
This proves (TDA.3); $\lfloor n/5\rfloor+1>n/5$ gives its
last assertion.
\end{proof}

\paragraph{Scope.}
This gives a degree bound only for an actual faithful projective
$A_n$-representation. An abstract quotient of a finite group
does not by itself furnish such a representation on a selected
subspace. The separate homogeneous central-product argument
must construct the needed multiplicity action, and the separate
conjugate-generation argument must supply any moving-rank bound.
No rank-nine source, polarization or detecting cycle follows
from (TDA.3) alone.
